\documentclass[fykos]{fksbase}
\setcounter{year}{27}

\geometry{ignoreheadfoot,a4paper,margin=1.5cm}
\fancyhf{}
\renewcommand{\headrulewidth}{0pt}

\usepackage{multicol}
\usepackage{colortbl}

\usepackage{minted}
\definecolor{mintedbg}{rgb}{0.90, 0.90, 1}
\newminted[texminted]{tex}{bgcolor=mintedbg, tabsize=2, xleftmargin=1mm, xrightmargin=1mm}
\newmintedfile[texfile]{tex}{bgcolor=mintedbg, tabsize=2, xleftmargin=1mm, xrightmargin=1mm,obeytabs=true}
\newmint{tex}{showspaces=false}
\newmintinline{tex}{showspaces=false}

\usepackage{mdframed}
\mdfsetup{backgroundcolor=blue!2, linecolor=blue!30, linewidth=3pt, leftline=false, rightline=false,
	innertopmargin=0, innerbottommargin=0, skipabove=15pt, skipbelow=15pt,}

\newcolumntype{T}{>{\columncolor{mintedbg}}m{0.3\textwidth}}
\newcolumntype{V}{>{\columncolor{blue!2}}m{0.3\textwidth}}
\newcolumntype{N}{>{\columncolor{white}}m{0.368\textwidth}}
\setlength\arrayrulewidth{0.5pt}\arrayrulecolor{blue!25}

\newenvironment{texcode}
	{\VerbatimEnvironment\hskip-10pt \minipage[t]{0.52\textwidth}\vspace{0pt}\begin{texminted}{tex}}
	{\end{texminted} \endminipage}
\newenvironment{texview}
	{\hskip 4pt \minipage[t]{0.50\textwidth}\vskip 4pt \vspace{0pt}\parindent=1em}
	{\vskip 6pt \endminipage}
\newenvironment{texexamplex}
	{\begin{mdframed}}
	{\vskip -2pt\end{mdframed}}
\newenvironment{texsource}
	{\VerbatimEnvironment\begin{mdframed}[backgroundcolor=mintedbg]\hskip-16pt\begin{texminted}{tex}}
	{\end{texminted}\hskip-40pt\end{mdframed}}
\newenvironment{texexample}
	{\VerbatimEnvironment\VerbatimOut[tabsize=2]{\jobname.tmp}}
	{\endVerbatimOut\begin{mdframed}\hskip-10pt\minipage[t]{0.52\textwidth}\vspace{0pt}\texfile{\jobname.tmp}\endminipage
		\hskip4pt\minipage[t]{0.50\textwidth}\vskip 4pt \vspace{0pt}\parindent=1em\input{\jobname.tmp}\vskip 6pt\endminipage\vskip-2pt\end{mdframed}}
\newenvironment{texexampletab}
	{\begin{mdframed}\hskip-10pt\begin{tabular}{TVN}}
	{\end{tabular}\end{mdframed}}
\begin{document}

\begin{center}
	\includegraphics{logo2.eps} \hspace{10pt}\raisebox{0.6cm}{\normalfont\sffamily\slshape\Huge Jak psát vzorová řešení}
	\vskip0.6cm
\end{center}

Tento text je určen zejména organizátorům FYKOSu/VÝFUKu a~má pomoci při
psaní vzorových řešení. Většinu zde popsaných příkazů však lze použít i~v~jiných
dokumentech a~usnadnit si tak práci hlavně při sazbě matematiky. V~tom případě
do hlavičky souboru přidejte balík \texttt{fkssugar} pomocí~{\texinline|\usepackage{fkssugar}|}.

\section{Překlad}

Překládáme pomocí \texttt{make}. Pokud nemůžete výchozí variantu použít,
vytvořte soubor \texttt{solution.tex} s~obsahem uvedeným níže. Změňte cestu
k~souboru úlohy a~přeložte příkazem \texttt{xelatex}. Obrázky nebudou do
souboru vloženy, protože v~některých případech by překlad tímto způsobem
kvůli chybějícím obrázkům selhal.

\begin{texsource}
\documentclass[fykos]{fkssolution} % nahradit jméno semináře
\setcounter{year}{27} % nahradit ročník
\renewcommand{\includegraphics}[2][]{\fbox{#2}}
\renewcommand\plotfig[4][h]{\begin{figure}[#1]\centering\fbox{#2}\caption{#3}\label{#4}\end{figure}}

\input{"fykos27/problems/problem1-7"} % nahradit soubor s úlohou

\begin{document}
	\problemsolutionsingle
	\makefooter
\end{document}
\end{texsource}

\section{Obecné zásady}

Řešení pište bez chyb! Kdo tam bude mít chyby, bude zmrskán!

Všechny texty týkající se FYKOSu jsou psány v~typografickém systému \TeX
(využívajíce {\LaTeX}ových maker). Řešení úloh pište v~libovolném
textovém editoru do souboru úlohy v~kódování UTF-8.

Výsledný soubor buď uložte
do repozitáře, nebo pošlete e-mailem TeXaři (v~tom případě je nutné je poslat
s~dostatečným předstihem před uzávěrkou). Do repozitáře nevkládejte soubory,
které nejdou přeložit. Pokud nemáte přístup do repozitáře a~chcete provést úpravu
v~řešení, stáhněte si aktuální verzi z~Astrid\footnote{\url{http://astrid.fykos.cz}},
upravte ji a~co nejdříve pošlete e-mailem TeXaři, aby se omezily konflikty úprav.

Text pište česky s~háčky a~čárkami v~první osobě množného čísla, asi po~60~znacích
(po~logických celcích) se snažte odřádkovávat. Více mezer a~(jednotlivé) odřádkování
nemá vliv na vizuální podobu výsledného textu. Nový odstavec vytvoříte prázdným
řádkem.

Text pište v~první osobě množného čísla, na konci svého řešení můžete uvést
hlavní chyby, kterých se řešitelé dopouštěli.


Velká většina příkazů se v~TeXu uvozuje zpětným lomítkem (\texinline|\|)
povinné argumenty se píší do složených závorek (\texinline|{}|), nepovinné
argumenty do závorek hranatých (\texinline|[]|). Například text do uvozovek
dáte pomocí~\texinline|\uv{}|.

Snažte se o~to, aby zdrojový kód byl co nejpřehlednější. Pokud nevíte, jaká je
přesná syntaxe určitého příkazu,
Google a~manuál {\LaTeX}u vám poradí. Poznámky, které se nemají objevit ve
výsledném textu, pište na řádek začínající \verb|%|.

\begin{texexample}
Lorem ipsum dolor sit amet, consectetur adipiscing.
\uv{Text v~uvozovkách} tint laoreet. Fusce varius
\textbf{tučně} vitae blandit. Cum sociis tetur
\textit{kurzívou}.

\emph{Zvýrazněný text} Venatibus et magnis,
nascetur ridiculus mus. Vestibulum eu neque
sem. Duis luctus vestibulum ultricies.

Nejsou tři tečky (výpustka) jako tři tečky\dots

\noindent\emph{Nápověda}\quad Takto vypíchneme
jedno slovo z odstavce (např. v zadání).
\end{texexample}

\section{Vlnky}

Vlnky (\texinline|~|) vkládáme tam, kde nechceme, aby došlo k~rozdělení na dva řádky. Patří:
\begin{compactitem}
	\item za jednopísmenná slova (a, i apod.),
	\item za zkratky,
	\item k veličinám (před nebo za, dle významu),
	\item k odkazům.
\end{compactitem}

\begin{texexample}
	Např.~Bob a~Bobek se pohybují rychlostí~$v$.
	Platí rovnice~\eqref{R27S1U2_rovnice-y}, přičemž
	$v$~je rychlost.
\end{texexample}

Pro doplnění vlnek na některá místa je možné použít skript \texttt{vlnka.sh} z~repozitáře \texttt{fks-scripts}.
Funguje pro češtinu a~slovenštinu a~vlnky doplňuje zejména na základě slovníku slov, u~kterých vlnka patří.
To však znamená, že může doplnit vlnky i~na místa, kde nejsou nutné, ale také tam, kde nepatří (do matematiky).
Naopak nedoplní vlnky všude. Proto je třeba po jeho spuštění soubor znovu projít a~smazat přebytečné a~přidat
nedoplněné vlnky. Pokud mažete vlnku z~matematiky, nevypisujte místo ní mezeru, aby se při dalším spuštění
skriptu vlnka opět nepřidala.

\section{Nadpisy}

V textu seriálu používejte nadpis \texinline|\subsection|, ve vzorových řešeních \texinline|\subsubsection|.

\section{Matematické vzorce}
\begin{compactitem}
	\item Inline matematiku píšeme mezi dolary (např.~\texinline|rychlost~$v_1 = w_0$|).
		V inline matematice nemají být zlomky zapsány jako zlomky, ale pomocí lomítka~\texinline|/|
		(např.~\texinline|rychlost $v=s/t$|).
	\item Jednořádkovou blokovou matematiku píšeme mezi \texinline|\eq{}| (ekvivalent \verb|$$|)
	a~víceřádkovou mezi \texinline|\eq[m]{}| (obdoba \texinline|\align|) a~před znaky, na něž chceme
	zarovnat, dáme \texinline|&|. Blokovou matematiku odsazujeme tabulátorem pro lepší čitelnost.
	\item Odkazy děláme pomocí \texinline|\lbl{ID}| na konkrétní řádek.
	\item Vzorec se chová jako část textu (čárku, tečku na konci řádku oddělujeme mezerou~\texinline|\,|).
	\item Ve vzorcích používejte normálně znaménka plus, mínus; znaménko krát není vhodné psát.
	\item Horní a~dolní index vytvoříte pomocí stříšky (\texinline|^|) nebo podtržítka~(\texinline|_|).
		Pokud je v~indexu více znaků, dejte je do složených závorek. Index, který
		vznikl z~nějakého slova, se píše stojatý~--~před podtržítko dejte zpětné lomítko~(\texinline|\_|).
		Výjimkou je tíhová síla $F_G$.
	\item Čárku nad veličinou pište pomocí anglického apostrofu (\texinline|'|).
	\item Pokud se v závorce vyskytuje zlomek, indexy apod., použijeme zvětšující se závorky
\end{compactitem}


\begin{texexample}
Rychlost míče~$v\_m=s/t$, ale rychlost míče
\eq{
	v\_m=\frac{s}{t} \,.
}
Tíhová síla~$F_G=ma$, ale gravitační síla
\eq{
	F\_g = \frac{G m_1 m_2}{r^2} \,.
}
Dále platí
\eq{
	y = ax + b \,, \lbl{R27S1U2_rovnice-y}
}
kde $a$~je směrnice a~$b$~posunutí.

Dle rovnice~\eqref{R27S1U2_rovnice-y} něco platí
a~něco zajímavého z toho určitě plyne. Řešme
soustavu rovnic
\eq[m]{
	x &= a+b+1 \,, \\
	y &= a-b \,.
}
Z~druhé vyjádříme~$b$ a~dosadíme do první.

Pokud je to třeba, použijte zvětšující se závorky
\eq{
	x = \left\lbrace \left[ \( \frac{x^2}{2} + 1 \)
	+ 1 \right] + 1 \right\rbrace \,.
}
\end{texexample}


\subsection{Příklady příkazů}

Hodit se vám budou následující příkazy:
\begin{multicols}{2}
\begin{compactitem}
	\item \verb|\frac{}{}| zlomek, do závorek napíšete čitatele a~jmenovatele,
	\item \verb|\cdot| tečka mezi dvěma zlomky,
	\item \verb|\sqrt{}| odmocnina,
	\item \verb|\alpha| malá alfa (ostatní řecká písmena obdobně),
	\item \verb|\Alpha| velká alfa,
	\item \verb|\int| integrál,
	\item \verb|\sum| součet,
	\item \verb|\sin{}| sinus,
	\item \verb|\cos{}| kosinus,
	\item \verb|\dot{}| tečka nad znakem (např. časová derivace),
	\item \verb|\ddot{}| dvě tečky (druhá časová derivace),
	\item \verb|\%| procento,
	\item \verb|\dots| tři tečky,
	\item \verb|\cdots| centrované tři tečky,
	\item \verb|\footnote{}| poznámka pod čarou.
\end{compactitem}
\end{multicols}

Další příkazy se týkají pouze FYKOSu (včetně Výfuku), jsou definovány ve FYKOSích
makrech
(balíček~\verb|fkssugar|).\nopagebreak
\begin{multicols}{2}
\begin{compactitem}
	\item \verb|\d| diferenciál,
	\item \verb|\der{y}{x}| první derivace $y$ podle $x$,
	\item \verb|\dder{}{}| druhá derivace,
	\item \verb|\pder{}{}| parciální derivace,
	\item \verb|\vect{}| vektor vysázený tučně,
	\item \verb|\bb{}| množiny (např. $\bb{R}$),
	\item \verb|\eu| Eulerovo číslo,
	\item \verb|\tg{}| tangens,
	\item \verb|\(| \verb|\)| levá a~pravá závorka,
	\item \verb|\ztoho| implikace,
	\item \verb|\const| konstanta,
	\item \verb|\dg| stupeň,
	\item \verb|\im| imaginární jednotka,
	\item \verb|\bod{A}| geometrický bod~$\bod{A}$,
	\item \verb|\C| stupně Celsia,
	\item \verb|\pi| kontanta $\approx "3.14"$,
	\item \verb|\oldpi| úhel, rovina, paralaxa.
\end{compactitem}
\end{multicols}

\begin{texexampletab}
	\texinline|$p\_{voda}$| & $p\_{voda}$ & dolní indexy, co nejsou veličina (kromě tíhového zrychlení) \\\hline
	\texinline|$\ztoho$| & $\ztoho$ & implikace \\\hline
%	\texinline|$\tsum_a^b$| & $\tsum_a^b$ & suma v řádku \\\hline
	\texinline|$\const$| & $\const$ & konstanta \\\hline
	\texinline|$\Re z, \Im z$| & $\Re z, \Im z$ & reálná a imaginární část \\\hline
	\texinline|$\eu^{\im x}$| & $\eu^{\im x}$ & Eulerovo číslo, imaginární jednotka \\\hline
	\texinline|$\int \d x$| & $\int \d x$ & integrál, diferenciál \\\hline
	\texinline|$\der{f}{x}, \dder{f}{x},|\ \texinline|\pder{f}{x}, \ppder{f}{x}$| & $\der{f}{x}, \dder{f}{x}, \pder{f}{x}, \ppder{f}{x}$ &\\\hline
	\texinline|$\Laplace, \Dalembert$| & $\Laplace, \Dalembert$ & \\\hline
	\texinline|$\rot, \Div, \grad$| & $\rot, \Div, \grad$ & \\\hline
	\texinline|$\prumer$| & $\prumer$ &\\\hline
	\texinline|$\eqdef,\doteq,\approx,\propto$| & $\eqdef, \doteq, \approx, \propto $ & definice, zaokrouhlení, aproximace, úměra\\\hline
	\texinline|$\qed$| & $\qed$ &\\\hline
	\texinline|$\vect{F}\_G$| & $\vect{F}\_G$ &\\\hline
	\texinline|$\bod{AB}$| & $\bod{AB}$ &\\\hline
	\texinline|$\bb{R}$| & $\bb{R}$ &\\\hline
	\texinline|$\op{dup}$| & $\op{dup}$ &\\\hline
	\texinline|$\pi$| & $ \pi$ &svisle konstanta $\approx "3.14"$\\\hline
	\texinline|$\ohm$| & $\ohm$ &\\\hline
\end{texexampletab}


\section{Čísla a~jednotky}

Chceme-li napsat desetinné číslo, příp. číslo s~jednotkami, použijeme uvozovkového
pomocníka. Ten je citlivý na mezeru před jednotkou. Pokud je jednotka součinem
několika jiných, píšeme tečku. Čísla do uvozovek píšeme s~desetinnou
tečkou (snadno pak naformátujeme i~výstup z~Gnuplotu). Exponent desítky napíšeme
pomocí \texinline|e|. Čísla oddělujeme do skupin po třech mezerou \texinline|\,|.
Chceme-li psát samostatně jednotku, využijeme příkaz \texinline|\jd{}|.

\begin{texexample}
Teplota~$"21.3 \C"$, \\
úhel je~$"123\dg"$, \\
$\sigma = "5.67e-8 J.s^{-1}.m^{-2}.K^{-4}"$, \\
$\varphi = "(0.768\,7\pm 0.000\,3) arcsec"$, \\
vzestup o~$"13.2 \%"$, \\
délka~$\lambda = "600 nm" = "0.6 \micro m"$, \\
rychlost v~$\jd{m.s^{-1}}$,\\
konstanta~$\pi \doteq "3.141\,592\,65"$,\\
konstanta~$\eu \doteq "2.718"$.
\end{texexample}


\section{Chemické vzorce}

Při sazbě různých chemických vzorců a~reakcí si můžeme ušetřit práci balíkem \texttt{mhchem}, používáme příkaz \texinline|\ce|.

\begin{texexample}
Vzorec látky je~\ce{NO3-}, půlka~\ce{1/2H2O},
izotop~\ce{^{227}_{90}Th+}. Chemická reakce
\eq{
	\ce{CO2 + C -> 2CO} \,. \lbl{R27S1U2_reakce}
}
Odkážeme se na rovnici~\eqref{R27S1U2_reakce}.
Můžeme i~nakreslit vazby,
např.:~\eq{
	\ce{A\bond{-}B\bond{=}C\bond{#}D} \,.
}
Objem vody je~$V_{\ce{H2O}}="5.1 m^3"$.
\end{texexample}


\section{Obrázky}

Obrázky jsou rozděleny do tří kategorií: grafy, datové/velké ilustrační a~drobné ilustrační (obtékané textem).

\tex|\plotfig{cesta k~souboru .tex}{popiska}{referenční ID} % Gnuplot term eps latex|
\tex|\fullfig{cesta k~souboru}{popiska}{referenční ID}|
\tex|\illfig{cesta k~souboru}{popiska}{referenční ID}{výška v~řádcích}|

Jediným povinným argumentem je cesta k~souboru (implicitně se prohledávají adresáře \verb|graphics|).
Ostatní argumenty lze též zadat jako prázdné tokenem \texinline|{}|. Chceme-li upravit velikost, popř.~jiný
argument velkého obrázku, napíšeme jej do nepovinného prvního argumentu.

\tex|\fullfig[width=0.6\textwidth]{cesta k~souboru}{popiska}{referenční ID}|

Obrázky (i~zdrojáky) k~úlohám se ukládají do adresáře \verb|problems/graphics|. Ostatní obrázky (seriál,
úvodníček) se ukládají k~dané sérii (\verb|batchB/graphics|).

\begin{figure}[h!]
\begin{texexamplex}
\begin{texcode}
\illfig{obrazek.eps}{Malý}{R27S1U2_maly}{4}
Obrázek~\ref{R27S1U2_maly} je malinký a~obtéká
ho text. Obrázek~\ref{R27S1U2_velky} je velký.

Lorem ipsum dolor sit amet, consectetur adipiscing
elit. Etiam neque urna, semper quis diam
pellentesque, scelerisque vehicula ante. Maecenas
et feugiat orci. Donec neque eros, eleifend sit amet
adipiscing vel, dictum at arcu. Pellentesque vitae
est ante.



\fullfig{obrazek.eps}{Velký obrázek}{R27S1U2_velky}





Quisque a~urna erat. Suspendisse in est sem. In
hac habitasse platea dictumst. Donec varius,
turpis a~tristique gravida, ligula nunc lobortis.
\end{texcode}
\begin{texview}
	\illfig{obrazek.eps}{Malý}{R27S1U2_maly}{6}
	
	Obrázek~\ref{R27S1U2_maly} je malinký a~obtéká ho text, kdežto
	obrázek~\ref{R27S1U2_velky} je velký.
	
	Lorem ipsum dolor sit amet, consectetur adipiscing elit. Etiam neque urna,
	semper quis diam pellentesque, scelerisque vehicula ante. Maecenas et feugiat
	orci. Donec neque eros, eleifend sit amet adipiscing vel, dictum at arcu.
	Pellentesque vitae est ante.
	
	\begin{center}
		\includegraphics{obrazek.eps}
		\caption{Velký obrázek}
		\label{R27S1U2_velky}
	\end{center}
	
	Quisque a~urna erat. Suspendisse in est sem. In hac habitasse platea dictumst.
	Donec varius, turpis a~tristique gravida, ligula nunc lobortis.
\end{texview}
\end{texexamplex}
\end{figure}

\section{Tabulky}

Nad tabulkou se sází vodorovná čára pomocí \texinline|\toprule|, \texinline|\midrule| odděluje
většinou záhlaví tabulky od zbytku a~pod tabulku patří vodorovná čára \texinline|\bottomrule|.

Popisek (\texinline|\caption{Popisek}|) se sází nad tabulku. Referenční ID (\texinline|\label{ID}|)
udáváme pro přehlednost ve tvaru např.~\verb|R27S1U2_id|, kde zaměníme ročník, sérii a~číslo
úlohy. Příkaz \texinline|\label| musí být vložen až po \texinline|\caption|. Na tabulku se
odkazujeme pomocí \texinline|\ref{ID}|.

V~záhlaví tabulky použijeme příkaz \texinline|\popit|, který funguje stejně jako \texinline|\popi|
(tedy vysází zlomek s~veličinou v~čitateli a~jednotkou ve jmenovateli), ale je pro použití
v~tabulkách, automaticky zarovnává popisek na střed sloupce. Pomocí nepovinných parametrů
je možné jej rozšířit na více sloupců a~jinak zarovnat, např.~pro dva sloupce a~zarovnání vlevo
použijeme \texinline|\popit[2][l]{t}{s} |.

Čísla v~tabulce by měla být zarovnána většinou podle desetinné čárky, popř.~vpravo. Záhlaví
zarovnejte doprostřed (použitím \texinline|\popit| nebo \texinline|\multicolumn|, pokud nechcete
sázet zlomek s~jednotkou).

\begin{table}[h!]
\begin{texexamplex}
\begin{texcode}
V~tabulce~\ref{R27S1U2_moje-tabulka} je mnoho dat.
Lorem ipsum dolor sit amet, imperdiet metus nunc,
eu faucibus nulla euismod ac. Fusce eu consequat mi.

\begin{table}[h!]
	\centering
	\caption{Moje tabulka}
	\label{R27S1U2_moje-tabulka}
	\begin{tabular}{rrrrrr}
		\toprule
		\multicolumn{1}{c}{$n$} & \popit{d}{cm} &
		\popit{|a'_1|}{m} & \popit{v}{m.s^{-2}} &
		\popit{t_2-t_1}{s} & $\sin{\alpha}$ \\
		\midrule
		1\,113 & 23 & 9 & 3,3 & $"2.1\pm 0.1"$ & 0,11 \\
		1\,233 & 22 & 8 & 4,5 & $"2.1\pm 0.1"$ & 0,17 \\
		1\,313 & 21 & 6 & 4,6 & $"2.1\pm 0.1"$ & 0,21 \\
		\bottomrule
	\end{tabular}
\end{table}

Integer enim nisl, sodales quis porta vel, tincidunt
a~elit. Fusce non ipsum egestas, egestas enim in,
facilisis enim. Vestibulum lacinia viverra est, at.
\end{texcode}
\begin{texview}
	V~tabulce~\ref{R27S1U2_moje-tabulka} je mnoho dat.
	Lorem ipsum dolor sit amet, imperdiet metus nunc,
	eu faucibus nulla euismod ac. Fusce eu consequat mi.
	
	
	\begin{center}
	\caption{Moje tabulka}
	\label{R27S1U2_moje-tabulka}
	\begin{tabular}{rrrrrr}
		\toprule
		\multicolumn{1}{c}{$n$} & \popit{d}{cm} &
		\popit{|a'_1|}{m} & \popit{v}{m.s^{-2}} &
		\popit{t_2-t_1}{s} & $\sin{\alpha}$ \\
		\midrule
		1\,113 & 23 & 9 & 3,3 & $"2.1\pm 0.1"$ & 0,11 \\
		1\,233 & 22 & 8 & 4,5 & $"2.1\pm 0.1"$ & 0,17 \\
		1\,313 & 21 & 6 & 4,6 & $"2.1\pm 0.1"$ & 0,21 \\
		\bottomrule
	\end{tabular}
	\end{center}
	
	\medskip
	
	Integer enim nisl, sodales quis porta vel, tincidunt a~elit. Fusce non ipsum egestas,
	egestas enim in, facilisis enim. Vestibulum lacinia viverra est, at.
\end{texview}
\end{texexamplex}
\end{table}

\section{Seznamy}

Seznam vytvoříte pomocí {\LaTeX}ového prostředí~\texinline|\compactitem| (nečíslovaný)
nebo~\texinline|\compactenum| (číslovaný, druh číslování můžete uvést v~nepovinné parametru),
každou položku označte~\texinline|\item|.

\begin{texexample}
Nákupní seznam s~písmenky:
\begin{compactenum}[a)]
	\item Lorem ipsum dolor sit amet, consectetur adi
		elit. Etiam neque urna, semper quis diam pelent,
		scelerisque vehicula ante.
	\item Maecenas et feugiat orci.
	\item Donec neque eros, eleifend sit amet.
\end{compactenum}

Nákupní seznam s~odrážkami:
\begin{compactitem}
	\item Nulla nec orci imperdiet, dapibus risus vit
		lacinia vestibulum, eleifend purus.
	\item Sed consequat mollis lacinia vestibulum.
\end{compactitem}
\end{texexample}



\section{Poznámky pod čarou}

Poznámky pod čarou se dělají pomocí příkazu~\texinline|\footnote{}|. Číslují se automaticky
a~vysází se v~dolní části strany. Pokud píšeme poznámku pod čarou uprostřed věty, píše se
ihned ke slovu, ke kterému patří. Na konci věty se píše za tečkou, pokud patří k~celé větě,
popř.~před tečkou, vztahuje-li se k~poslednímu slovu věty.

\begin{texexample}
Lorem ipsum\footnote{Poznámka ke slovu} dolor
sit amet, consectetur adipis cing elit.

Vestibulum nec augue.\footnote{Poznámka
k~celé větě} Pel lentesque convallis ac felis tinci
dunt rhoncus. Cras eleifend tortor sit amet.
\end{texexample}


\section{Ostatní}

Na webové odkazy použijte makro \texinline|\url{adresa}|.

Odkazy do textu vložíte pomocí maker~\texinline|\ref{referenční ID}| (obrázky, tabulky)
a~\texinline|\eqref{referenční ID}| (pro rovnice, číslo bude v~závorkách). Neodkazujte
se takto na rovnice, tabulky a~obrázky mimo řešení příkladu. Číslování je v~různých
souborech (brožurka, samotné řešení, ročenka) různé, takže není jednoznačné.


\section{Vzorový soubor}

Následující soubor je minimálním příkladem souboru, ve kterém bude fungovat téměř
vše, co je v~tomto dokumentu uvedeno (kromě obrázků, ty je nutno vkládat běžným
způsobem). Můžete tak zde popsané postupy (včetně FYKOSích specialit jako
např.~snadné psaní čísel s jednotkami pomocí uvozovek) použít i ve všech ostatních
dokumentech. Překládáme příkazem \texttt{xelatex}.

\begin{texsource}
\documentclass[10pt,a4paper]{article} % typ dokumentu, velikost písma a strany
\usepackage{amssymb}
\usepackage[intlimits]{amsmath}
\usepackage[no-math]{fontspec}
\usepackage{xltxtra,polyglossia}
\setmainlanguage{czech} % dokument bude v češtině
\mathcode`\,="013B % používá se desetinná čárka
\usepackage{fkssugar} % FYKOS
\usepackage[margin=2cm]{geometry} % okraje stránky
\usepackage{color,graphicx,graphics} % barvičky, obrázky
\usepackage{booktabs,paralist,url} % tabulky, seznamy, odkazy
\usepackage[version=3]{mhchem} % chemie
\newcommand\uv[1]{„#1“} % uvozovky

\author{Tomáš Pikálek}
\title{Můj krásný dokument}
\begin{document}
	Můj první dokument v~\LaTeX u.
\end{document}
\end{texsource}

\end{document}